\documentclass[11pt,letterpaper]{article}

\usepackage[utf8]{inputenc}
\usepackage{amsmath}
\usepackage{amssymb}
\usepackage{graphicx}
\usepackage{booktabs}
\usepackage{hyperref}
\usepackage{float}
\usepackage{geometry}
\usepackage{xcolor}
\usepackage{listings}

\geometry{margin=0.75in}

\title{Avance 6: Conclusiones Clave \\
\large Viabilidad de Implementación y Estrategia de Despliegue en Producción \\
VisionOps -- Reconocimiento de Acciones Humanas en Manufactura}

\author{
  Landy Haydee Schlebach Osorio \\
  Carlos Pano Hernández \\
  Carlos Fernando Del Castillo Rey \\[0.5em]
  \small Tec de Monterrey, Campus Estado de México -- Equipo MNA-V \\
  \small Asesor académico: Dr. Gerardo Camacho \\
  \small Patrocinador industrial: Alignity IQ Edge, LLC
}

\date{Junio 2026}

\begin{document}

\maketitle

\begin{abstract}
Este avance final (Avance 6) evalúa la viabilidad de implementación del modelo V-JEPA 2 en producción. Se realiza análisis exhaustivo de criterios de éxito, rendimiento del modelo, recomendaciones clave y selección de plataforma cloud. Tras evaluación de AWS, Azure, GCP e IBM Watson, se recomienda AWS como plataforma prioritaria por su superioridad en escalabilidad, costo y servicios ML-específicos. Se proporciona arquitectura de producción, procedimientos de implementación y tareas accionables para stakeholders.
\end{abstract}

\section{Ejecutivo}

\subsection{Verdict de Implementación}

\textbf{RECOMENDACIÓN: IMPLEMENTAR EN PRODUCCIÓN} [APROBADO]

El modelo V-JEPA 2 cumple con todos los criterios de éxito establecidos:
\begin{itemize}
  \item Accuracy 78.54\% (exceeds 75\% target)
  \item AUC-ROC 0.9573 (excellent discrimination)
  \item Latencia <5ms (meets real-time requirement)
  \item Confiabilidad de calibración: 0.85 (good)
  \item No requiere revisión de fases anteriores
\end{itemize}

\textbf{Plataforma Recomendada: AWS} (SageMaker + Lambda + RDS)

\section{Evaluación de Viabilidad}

\subsection{Criterios de Éxito Establecidos}

\begin{table}[H]
\centering
\caption{Validación contra Criterios de Éxito Predefinidos}
\begin{tabular}{|l|c|c|c|}
\hline
\textbf{Criterio} & \textbf{Meta} & \textbf{Logrado} & \textbf{Status} \\
\hline
Accuracy en validación & $\geq 75\%$ & 78.54\% & \textcolor{green}{OK} \\
\hline
Macro F1-Score & $\geq 0.70$ & 0.7063 & \textcolor{green}{OK} \\
\hline
AUC-ROC promedio & $\geq 0.90$ & 0.9573 & \textcolor{green}{OK} \\
\hline
Latencia inferencia & $< 10$ ms & $< 5$ ms & \textcolor{green}{OK} \\
\hline
Confianza calibrada & $> 0.80$ & 0.85 & \textcolor{green}{OK} \\
\hline
Cobertura de clases & 12/12 & 12/12 & \textcolor{green}{OK} \\
\hline
Estabilidad del modelo & $\pm 2\%$ & $<1\%$ & \textcolor{green}{OK} \\
\hline
\end{tabular}
\end{table}

\subsection{Análisis de Readiness}

\textbf{Fortalezas del Modelo:}
\begin{itemize}
  \item Discriminación excelente (AUC 0.9573)
  \item Confianza bien calibrada (0.85 media)
  \item Latencia compatible con tiempo real (<5ms)
  \item Backbone pre-entrenado (V-JEPA 2) robusto
  \item Embeddings congelados = reproducibilidad garantizada
  \item Simple architecture (Logistic Regression) = fácil mantenimiento
\end{itemize}

\textbf{Áreas de Mejora Identificadas:}
\begin{itemize}
  \item Clases minoritarias (n<50) tienen accuracy variable (50-70\%)
  \item Accuracy global 78.54\% deja margen para 21.46\% de mejora
  \item Confusiones entre clases similares (picking left/front)
  \item Dataset actual limitado a 3,425 ejemplos
\end{itemize}

\section{Respuesta a Preguntas Clave}

\subsection{¿Rendimiento suficiente para producción?}

\textbf{SÍ.} El modelo alcanza 78.54\% accuracy superando la meta de 75\%. En contexto industrial:
\begin{itemize}
  \item AUC 0.9573 indica excelente separabilidad de clases
  \item Confianza 0.85 permite filtro de predicciones inciertas
  \item Latencia <5ms permite monitoreo en tiempo real
  \item Matriz confusión muestra errores predecibles (no aleatorios)
  \item Recomendación: Usar threshold confianza 0.75 en producción
\end{itemize}

\subsection{¿Existe margen para mejorar?}

\textbf{SÍ, significativamente.} Oportunidades identificadas:

\begin{enumerate}
  \item \textbf{Aumentar datos:} De 3,425 a 10,000+ ejemplos (SMOTE para clases raras)
  \item \textbf{Fine-tuning selectivo:} Descongelar últimas capas de V-JEPA 2 con datos reales
  \item \textbf{Ensambles adaptados:} Entrenar meta-learner con datos reales (no validación)
  \item \textbf{Augmentación:} Rotaciones, zoom, oclusión en videos
  \item \textbf{Nuevas arquitecturas:} Vision Transformers (ViT) como alternativa
\end{enumerate}

Estimación: Mejora potencial hasta 85-88\% con estrategias anteriores.

\subsection{Recomendaciones Clave de Implementación}

\begin{table}[H]
\centering
\caption{Recomendaciones Prioritarias por Fase}
\begin{tabular}{|l|p{4cm}|l|}
\hline
\textbf{Fase} & \textbf{Recomendación} & \textbf{Prioridad} \\
\hline
\textbf{Inmediata} & Desplegar V-JEPA 2 en AWS SageMaker & CRÍTICA \\
\hline
& Configurar monitoring y alertas & CRÍTICA \\
\hline
& Implementar threshold confianza 0.75 & ALTA \\
\hline
\textbf{Corto plazo} & Recolectar datos reales de planta & ALTA \\
(1-3 meses)
& Validación HITL con operadores & ALTA \\
\hline
& Fine-tuning selectivo con datos reales & MEDIA \\
\hline
\textbf{Mediano plazo} & Aumentar dataset a 10,000+ ejemplos & MEDIA \\
(3-6 meses)
& Explorar ensambles meta-learner adaptados & MEDIA \\
\hline
\textbf{Largo plazo} & Evaluación de Vision Transformers & BAJA \\
(6+ meses)
& Análisis de interpretabilidad (SHAP/LIME) & BAJA \\
\hline
\end{tabular}
\end{table}

\subsection{Tareas Accionables para Stakeholders}

\subsubsection{Rol: Data Engineering}

\begin{itemize}
  \item [*] Crear pipeline de ingesta de datos en tiempo real desde YOLO detector
  \item [*] Configurar almacenamiento en S3 con política de retención (90 días)
  \item [*] Validar calidad de embeddings (V-JEPA y DINOv2) en producción
  \item [*] Implementar logging de predicciones y confianza en RDS
  \item [*] Automatizar reentrenamiento semanal con datos nuevos
\end{itemize}

\subsubsection{Rol: DevOps / Cloud Architecture}

\begin{itemize}
  \item [*] Configurar VPC en AWS con seguridad TLS/SSL
  \item [*] Desplegar modelo en SageMaker Endpoint (multi-AZ)
  \item [*] Crear Lambda function para inference (cold start <500ms)
  \item [*] Configurar autoscaling basado en throughput (100-500 req/s)
  \item [*] Implementar CloudWatch dashboards para monitoreo
\end{itemize}

\subsubsection{Rol: ML Operations}

\begin{itemize}
  \item [*] Establecer baseline de drift detection (accuracy <75\% = alert)
  \item [*] Definir SLOs: 99.9\% uptime, <10ms latency p95
  \item [*] Crear procedimiento de rollback en 5 minutos
  \item [*] Validar modelo monthly contra nuevos datos
  \item [*] Documentar decisiones en MLflow
\end{itemize}

\subsubsection{Rol: Production Floor / Manufacturing}

\begin{itemize}
  \item [*] Entrenar operadores en interpretación de alertas (confianza < 0.75)
  \item [*] Definir acciones cuando modelo predice con baja confianza
  \item [*] Recolectar feedback de operadores (casos borderline)
  \item [*] Registrar mejoras observadas en OEE (Overall Equipment Effectiveness)
  \item [*] Participar en validación HITL (human-in-the-loop)
\end{itemize}

\section{Evaluación de Plataformas Cloud}

\subsection{Comparativa AWS vs Azure vs GCP vs IBM Watson}

\begin{table}[H]
\centering
\caption{Comparativa Multidimensional de Plataformas Cloud}
\begin{tabular}{|l|cccc|}
\hline
\textbf{Aspecto} & \textbf{AWS} & \textbf{Azure} & \textbf{GCP} & \textbf{IBM Watson} \\
\hline
\textbf{Escalabilidad} & 5/5 & 4.5/5 & 4.5/5 & 3.5/5 \\
\hline
\textbf{Servicios ML} & 5/5 & 4/5 & 5/5 & 3/5 \\
\hline
\textbf{Facilidad de Uso} & 4/5 & 4.5/5 & 4/5 & 3.5/5 \\
\hline
\textbf{Costo} & 4/5 & 3.5/5 & 4.5/5 & 3/5 \\
\hline
\textbf{Integración Datos} & 5/5 & 4.5/5 & 4.5/5 & 3/5 \\
\hline
\textbf{Soporte Técnico} & 5/5 & 4.5/5 & 4/5 & 3.5/5 \\
\hline
\textbf{Comunidad} & 5/5 & 4.5/5 & 4.5/5 & 2.5/5 \\
\hline
\textbf{Experiencia del Equipo} & 4/5 & 3/5 & 3/5 & 2/5 \\
\hline
\hline
\textbf{SCORE TOTAL} & \textbf{37/40} & \textbf{32.5/40} & \textbf{33.5/40} & \textbf{24.5/40} \\
\hline
\end{tabular}
\end{table}

\subsection{Análisis Detallado por Proveedor}

\subsubsection{AWS (Amazon Web Services)}

\textbf{Fortalezas:}
\begin{itemize}
  \item SageMaker: platform ML nativa con features avanzadas (model registry, feature store)
  \item Lambda: serverless para inference con auto-scaling
  \item RDS: base datos relacional managed para predicciones
  \item S3: almacenamiento escalable con versioning
  \item CloudWatch: monitoring y logging integrado
  \item Experiencia del equipo: Alto (múltiples proyectos previos)
  \item Costo: Competitivo con descuentos por volumen
  \item Comunidad: Gigantesca, recursos abundantes
\end{itemize}

\textbf{Debilidades:}
\begin{itemize}
  \item Curva de aprendizaje para servicios avanzados
  \item Precio inicial puede parecer alto
  \item Requiere gestión de IAM roles/policies
\end{itemize}

\textbf{Arquitectura Propuesta (AWS):}
\begin{verbatim}
[YOLO Detector] -> [Lambda] -> [SageMaker Endpoint] -> [RDS]
                      |
                      +-> [CloudWatch] -> [Alerts]
\end{verbatim}

\subsubsection{Azure (Microsoft Azure)}

\textbf{Fortalezas:}
\begin{itemize}
  \item Azure ML: plataforma integrada con VSCode
  \item Integración con ecosistema Microsoft (Teams, Office)
  \item CosmosDB: base datos NoSQL con replicación global
  \item Precio competitivo para empresas con licencias Microsoft
\end{itemize}

\textbf{Debilidades:}
\begin{itemize}
  \item Servicios menos especializados que AWS/GCP
  \item Documentación fragmentada entre servicios
  \item Experiencia del equipo: Baja
\end{itemize}

\subsubsection{Google Cloud Platform (GCP)}

\textbf{Fortalezas:}
\begin{itemize}
  \item Vertex AI: plataforma unificada para ML
  \item BigQuery: almacenamiento + análisis masivo
  \item TensorFlow: framework nativo con soporte completo
  \item Infraestructura de datos robusta
\end{itemize}

\textbf{Debilidades:}
\begin{itemize}
  \item Menos maduro que AWS en algunos servicios
  \item Comunidad más pequeña
  \item Experiencia del equipo: Media
\end{itemize}

\subsubsection{IBM Watson}

\textbf{Fortalezas:}
\begin{itemize}
  \item Especializado en enterprise AI
  \item AutoML capabilities
\end{itemize}

\textbf{Debilidades:}
\begin{itemize}
  \item Precio muy elevado
  \item Comunidad limitada
  \item Menos servicios que competidores
  \item Experiencia del equipo: Nula
  \item Overkill para este proyecto
\end{itemize}

\subsection{Justificación de Selección: AWS}

\textbf{RECOMENDACIÓN FINAL: AWS}

AWS es la plataforma más apropiada por las siguientes razones:

\begin{enumerate}
  \item \textbf{Escalabilidad:} SageMaker soporta millones de predicciones/día con auto-scaling
  \item \textbf{Costo-beneficio:} Pricing por inferencia + Lambda serverless = economía
  \item \textbf{Servicios integrados:} S3 + RDS + Lambda + CloudWatch = arquitectura completa
  \item \textbf{Experiencia del equipo:} Equipo ha trabajado con AWS previamente
  \item \textbf{Comunidad:} Recursos, documentation y soporte abundante
  \item \textbf{Seguridad:} Compliance con GDPR, HIPAA, ISO 27001
  \item \textbf{Roadmap:} Actualizaciones constantes de servicios ML
\end{enumerate}

\textbf{Score Final: 37/40 (92.5\%)}

\section{Arquitectura de Producción Propuesta}

\subsection{Componentes Principales}

\begin{table}[H]
\centering
\caption{Componentes de la Arquitectura de Producción}
\begin{tabular}{|l|l|p{4cm}|}
\hline
\textbf{Componente} & \textbf{Servicio AWS} & \textbf{Función} \\
\hline
Ingesta de datos & API Gateway + Lambda & Recibir embeddings desde YOLO \\
\hline
Almacenamiento temporal & S3 & Guardar embeddings y predicciones \\
\hline
Modelo ML & SageMaker Endpoint & Servir predicciones (inference) \\
\hline
Base de datos & RDS (PostgreSQL) & Persistencia de predicciones \\
\hline
Orquestación & Lambda + EventBridge & Trigger reentrenamiento automático \\
\hline
Monitoring & CloudWatch & Alertas y dashboards \\
\hline
Autoscaling & Auto Scaling Groups & Escalar Lambda/Endpoint según demanda \\
\hline
\end{tabular}
\end{table}

\subsection{Flujo de Datos en Producción}

\begin{verbatim}
Planta (YOLO)
    |
    v
[API Gateway]
    |
    v
[Lambda Function] (validar + pre-procesar)
    |
    +---> [SageMaker Endpoint]
    |          |
    |          v (predicción)
    |     [Logistic Regression]
    |          |
    |          v
    +---> [RDS] (almacenar predicción + confianza)
    |
    +---> [CloudWatch] (monitoreo)
    |
    +---> [SNS] (alertas si confianza < 0.75)
    |
    v
[Dashboard] (operadores planta)
\end{verbatim}

\subsection{Especificaciones Técnicas}

\textbf{SageMaker Endpoint:}
\begin{itemize}
  \item Instancia: ml.m5.xlarge (4 vCPU, 16 GB RAM)
  \item Modelo: V-JEPA 2 + Logistic Regression (serializado en joblib)
  \item Timeout: 30 segundos
  \item Concurrencia: 100 instancias (escala automática)
  \item Latencia esperada: 2-5ms por predicción
\end{itemize}

\textbf{Lambda Function:}
\begin{itemize}
  \item Runtime: Python 3.11
  \item Memory: 512 MB
  \item Timeout: 60 segundos
  \item Triggers: API Gateway (síncrono), EventBridge (asíncrono)
  \item Concurrencia reservada: 1000
\end{itemize}

\textbf{RDS PostgreSQL:}
\begin{itemize}
  \item Instancia: db.t3.medium (2 vCPU, 4 GB RAM)
  \item Storage: 100 GB gp3 (escalable)
  \item Backup: Diario, retención 30 días
  \item Multi-AZ: Sí (alta disponibilidad)
  \item Replicación de lectura: 2 réplicas para análisis
\end{itemize}

\section{Procedimientos de Implementación}

\subsection{Cronograma de Despliegue (12 semanas)}

\begin{table}[H]
\centering
\caption{Timeline de Implementación en Producción}
\begin{tabular}{|l|p{3cm}|l|}
\hline
\textbf{Semana} & \textbf{Hito} & \textbf{Deliverable} \\
\hline
1-2 & Setup AWS + IaC & VPC, IAM, S3, RDS configurados \\
\hline
3-4 & SageMaker Endpoint & Modelo desplegado, test unitarios pass \\
\hline
5 & Lambda Functions & API Gateway integrada, stress testing \\
\hline
6 & Staging Deployment & Ambiente pre-prod con datos reales \\
\hline
7 & Validación HITL & Operadores validan 500 predicciones \\
\hline
8 & Monitoring Setup & CloudWatch, alertas, dashboards \\
\hline
9 & Training Operadores & Workshops para planta \\
\hline
10 & Canary Deployment & 10\% tráfico a nuevo modelo \\
\hline
11 & Blue-Green Deployment & 100\% migración (rollback disponible) \\
\hline
12 & Post-Deployment & Validación, ajustes, documentación final \\
\hline
\end{tabular}
\end{table}

\subsection{Criterios de Éxito Fase Deployment}

\begin{itemize}
  \item SLO met: 99.9\% uptime en staging 1 semana
  \item Latencia: p95 < 10ms, p99 < 50ms
  \item Accuracy en staging: $\geq 78\%$ (vs 78.54\% en validación)
  \item Falsos positivos: $< 2\%$
  \item Operadores entrenados: 100\%
  \item Documentación completa: SOP, runbooks, disaster recovery
\end{itemize}

\section{Gestión de Riesgos}

\subsection{Riesgos Identificados y Mitigación}

\begin{table}[H]
\centering
\caption{Matriz de Riesgos y Mitigación}
\begin{tabular}{|l|l|l|l|}
\hline
\textbf{Riesgo} & \textbf{Probabilidad} & \textbf{Impacto} & \textbf{Mitigación} \\
\hline
Drift del modelo & Media & Alto & Monitoring mensual \\
en nuevos datos & & & Fine-tuning trimestral \\
\hline
Latencia >10ms & Baja & Alto & Cacheing, SageMaker Async \\
en prod & & & Replicación de lectura \\
\hline
Outage SageMaker & Muy baja & Crítico & Multi-AZ, fallback local \\
\hline
Data quality issue & Media & Medio & Validación de embeddings \\
\hline
Security breach & Baja & Crítico & TLS/SSL, VPC, IAM roles \\
\hline
\end{tabular}
\end{table}

\section{Conclusión y Recomendación Final}

\subsection{Veredicto de Implementación}

\textbf{RECOMENDACIÓN: IMPLEMENTAR INMEDIATAMENTE}

El modelo V-JEPA 2 cumple todos los criterios de éxito y está listo para despliegue en producción. Los resultados (78.54\% accuracy, 0.9573 AUC) demuestran viabilidad técnica. AWS es la plataforma seleccionada como infraestructura cloud debido a su escalabilidad, costo-efectividad y servicios especializados.

\subsection{Próximos Pasos Inmediatos}

\begin{enumerate}
  \item Aprobación del documento de implementación por stakeholders
  \item Asignación de equipo (DevOps, Data Engineer, ML Ops)
  \item Inicio de setup AWS (Semana 1)
  \item Planificación de reentrenamiento continuo
  \item Comunicación a operadores de planta
\end{enumerate}

\subsection{Sostenibilidad a Largo Plazo}

El modelo es sostenible en producción con:
\begin{itemize}
  \item Reentrenamiento mensual con datos nuevos
  \item Monitoreo continuo de drift y performance
  \item Rotación de on-call engineers para soporte 24/7
  \item Budget anual estimado: \$50-80K USD (AWS + equipo)
  \item Mejora esperada: De 78.54\% a 85-88\% en 6 meses
\end{itemize}

\end{document}
